\documentclass[12pt]{article}
\usepackage[T1]{fontenc}
\usepackage[utf8]{inputenc}
\usepackage{lmodern}
\usepackage{textcomp}
\usepackage{enumitem}
\usepackage{geometry}
\geometry{margin=1in}
\begin{document}
\begin{center}
Answer Key - Economics - Undergraduate Level Assessment 25 \
\small (Answers are ordered exactly as in the question paper.)
\end{center}

\section*{Multiple Choice}
\begin{enumerate}[label=\arabic*.]
\item \textbf{Answer:} A
\\ \textit{Options:} A) Establishing minimum standards, providing public services, setting forth laws; B) Funding education, providing healthcare, controlling inflation; C) Printing money, establishing trade agreements, military defense
\\ \textit{Note:} The correct answer highlights the government's essential functions in regulating economic activity by establishing minimum standards to protect consumers and workers, providing public services that benefit society as a whole, and enacting laws that ensure fair competition and stability in the marketplace.
\item \textbf{Answer:} C
\\ \textit{Options:} A) Purchasing power increases and more goods can be purchased with the same amount of money.; B) Purchasing power remains the same.; C) Purchasing power decreases and less goods can be purchased with the same amount of money.; D) Purchasing power fluctuates unpredictably.
\\ \textit{Note:} The correct answer is based on the key concept that inflation reduces the value of money over time, meaning that as prices for goods and services rise, consumers are able to buy fewer items with the same amount of money, thus decreasing purchasing power.
\item \textbf{Answer:} A
\\ \textit{Options:} A) supply falls, increasing the price and increasing the quantity.; B) demand rises, decreasing the price and increasing the quantity.; C) demand rises, increasing the price and decreasing the quantity.; D) demand falls, decreasing the price and increasing the quantity.; E) supply rises, decreasing the price and increasing the quantity.
\\ \textit{Note:} The correct answer is based on the principle that when the price of a substitute good, such as polyester pants, rises, consumers are likely to switch their demand to denim jeans, leading to an increase in demand for jeans, which can cause the price to rise and incentivize suppliers to reduce their quantity supplied, ultimately resulting in a higher equilibrium price and quantity.
\item \textbf{Answer:} D
\\ \textit{Options:} A) efficiency is achieved; B) there is no opportunity cost; C) utility is maximized; D) resources are not specialized; E) there is economic growth
\\ \textit{Note:} A production possibility frontier will be a straight line when resources are not specialized because this indicates that the opportunity cost of reallocating resources between the production of two goods remains constant, reflecting a uniform trade-off.
\item \textbf{Answer:} A
\\ \textit{Options:} A) Higher interest rates that result from borrowing to conduct contractionary fiscal policy; B) Lower interest rates that result from borrowing to conduct expansionary monetary policy; C) Lower interest rates due to borrowing to conduct contractionary fiscal policy; D) Lower interest rates resulting from reduced borrowing due to expansionary fiscal policy; E) Lower interest rates due to increased borrowing to conduct expansionary monetary policy
\\ \textit{Note:} The "crowding-out" effect occurs when increased government borrowing to finance contractionary fiscal policy raises interest rates, which in turn discourages private investment due to the higher cost of borrowing.
\end{enumerate}

\section*{True / False}
\begin{enumerate}[label=\arabic*.]
\setcounter{enumi}{5}
\item \textbf{Answer:} True
\\ \textit{Options:} A) True; B) False
\\ \textit{Note:} In monopolistically competitive industries, the presence of profits attracts new firms to enter the market, leading to an increase in the number of firms.
\item \textbf{Answer:} True
\\ \textit{Options:} A) True; B) False
\\ \textit{Note:} The GDP Deflator measures inflation by comparing the nominal GDP, which uses current year prices and quantities, to the real GDP, which is adjusted for inflation using base year prices, thus incorporating current year quantities in its calculations.
\item \textbf{Answer:} False
\\ \textit{Options:} A) True; B) False
\\ \textit{Note:} The statement is false because the velocity of money is calculated as the ratio of nominal GDP to the money supply, which in this case is 3 (i.e., $600 billion / $200 billion), not 6.
\item \textbf{Answer:} True
\\ \textit{Options:} A) True; B) False
\\ \textit{Note:} A leftward shift in supply decreases the quantity available in the market, while a rightward shift in demand increases the desire for the product, leading to upward pressure on the equilibrium price of corn.
\item \textbf{Answer:} False
\\ \textit{Options:} A) True; B) False
\\ \textit{Note:} The statement is false because, in the long run, firms have the flexibility to adjust all factors of production, including labor, capital, and land, to optimize their production processes and respond to changes in market conditions.
\end{enumerate}

\section*{Long Form Response}
\begin{enumerate}[label=\arabic*.]
\setcounter{enumi}{10}
\item \textbf{Answer:} To calculate the annual depreciation for a machine costing $1,000,000 with a 10-year lifespan and a $100,000 scrap value, one can use the straight-line depreciation method. This involves subtracting the scrap value from the initial cost and then dividing by the useful life of the asset. Specifically, the calculation is ($1,000,000 - $100,000) / 10 years, resulting in an annual depreciation expense of $90,000. This depreciation expense impacts financial statements by reducing taxable income on the income statement and decreasing the asset value on the balance sheet, which reflects the machine's consumption over time. Therefore, the annual depreciation is $90,000, not \$70,000, indicating a need for clarity in the provided key concept.
\\ \textit{Note:} correct because it accurately applies the straight-line depreciation formula by accounting for the initial cost, scrap value, and useful life of the asset, resulting in a clear annual depreciation expense that directly affects taxable income on financial statements.
\item \textbf{Answer:} The use of Ordinary Least Squares (OLS) in the presence of autocorrelation leads to two significant implications: biased standard errors and unreliable forecasts. Autocorrelation violates the OLS assumption of independent errors, causing the estimated standard errors to be underestimated. This underestimation results in misleading inference, as hypothesis tests may incorrectly indicate statistical significance. Furthermore, forecasts derived from an OLS model with autocorrelation may be systematically biased, leading to poor predictions in economic modeling. Consequently, addressing autocorrelation is crucial for ensuring the reliability of both standard errors and forecasts in econometric analyses.
\\ \textit{Note:} correct because autocorrelation in OLS violates the assumption of independent errors, leading to biased standard errors that misrepresent statistical significance and unreliable forecasts due to the systematic nature of the error term.
\end{enumerate}

\vfill
\noindent\textit{End of Answer Key}
\end{document}